
\section{appendix}

\subsection{Analysis and Discussion}
Although AgentCAT demonstrates reliability in dynamic interaction and ability assessment, as an LLMs-based framework, its system overhead and result volatility are critical factors that must be considered in practical deployment and research.Within the multi-agent collaborative framework of AgentCAT, the primary computational overhead originates from the token consumption when the examinee agent generates long-text rationales and the selection agent conducts planning and retrieval. Taking the GLM-4-FlashX model as an example, according to the API call log statistics in our authentic testing environment, the average token overhead per API call within a continuously iterating long context is approximately 1200. For a complete adaptive testing sequence of length $T=20$, the total token consumption for the dual-end simulated interaction (including exercise retrieval and response simulation) of a single examinee is approximately 48,000. Under the current pricing standards of foundational large models, completing a full online dynamic assessment simulation for one examinee costs only about 0.005 RMB; even for the large-scale group simulation experiment with 30 concurrent users conducted in this study, the total inference cost is highly stably controlled within 0.15 RMB.

AgentCAT also exhibited certain failure cases during the experiments:~(1) Format Alignment Failure and Semantic Hallucination: In a minority of cases, the selection agent failed to correctly output the knowledge concept and bucket tags, or produced logical hallucinations. For instance, it might determine in its internal chain-of-thought that the examinee's ability is insufficient, yet select a high-difficulty feature bucket in its final output.~(2) Polarization and Local Oscillation in Ability Estimation: The experiments observed that in certain simulation trajectories, the system's ability estimation did not converge smoothly as expected, but instead exhibited severe overestimation or underestimation phenomena. This phenomenon likely stems from the distributional discrepancy between the simulated responses of large models and authentic human behaviors regarding the ``guessing'' and    ``slipping'' mechanisms. When the examinee agent, due to model hallucinations or over-reasoning, unexpectedly outputs the correct answer on a severely out-of-scope, highly difficult exercise, or provides an incorrect answer on a fundamental, error-prone exercise due to a logical lapse, the supervisor receives a massive penalty or reward gradient.


\subsection{Prompt}
\begin{figure}[htbp]
\centering
\begin{tcolorbox}[
    enhanced,
    width=\columnwidth, % 显式指定宽度填满双栏
    colback=gray!2!white,
    colframe=gray!75!black,
    title=\textbf{Prompt Template for Selection Agent},
    fontupper=\footnotesize\ttfamily,
    halign=flush left,
    arc=2mm,
    boxrule=0.5pt,
    left=2mm,
    right=2mm,
    top=1mm,
    bottom=1mm
    % 删除了 breakable 参数
]

% 第一部分：系统指令
\textbf{[System Instruction]} \\
You are a question selection agent, familiar with computerized adaptive testing principles, able to select appropriate questions based on student answering situations. \\
Please use english for all explanatory text; but must strictly use the english tags and bucket names given below for structured output: \\
1) Output tags: 'Knowledge:' and 'Bucket:' \\
2) Bucket names (choose one): \textcolor{blue}{\{bucket\_options\}} \\
3) Please copy the knowledge point name exactly from the 'Available Buckets Summary' below. \\
Final output must only include the two lines below.

\vspace{0.4em}
\tcbline
\vspace{0.4em}

% 第二部分：动态数据注入
\textbf{[Student Profile]} \\
\textcolor{blue}{\{student\_profile\_data\}} 

\vspace{0.5em}
\textbf{[Recent Answer Result]} \\
- Knowledge Point: \textcolor{blue}{\{current\_kc\}} \\
- KC Reliability: \textcolor{blue}{\{kc\_reliability\}} \\
- Answer: \textcolor{blue}{\{student\_answer\}} \\
- Score: \textcolor{blue}{\{score\}} \\
- Corrective Feedback: \textcolor{blue}{\{feedback\}} \\
- Knowledge Dependencies: \textcolor{blue}{\{dependencies\_str\}}

\vspace{0.8em}
\textbf{[Available Buckets Summary]} \\
\textcolor{blue}{\{stats\_text\}}

\vspace{0.8em}
\tcbline
\vspace{0.8em}

% 第三部分：具体执行要求
\textbf{[Instructions]} \\
Select the most appropriate knowledge point and bucket based on student ability, knowledge proficiency, dependency relationships, and availability. 
If KC Reliability is low ($<0.35$), rely less on the last predicted KC and prioritize ability/difficulty matching. Do not output specific exercise IDs.

\vspace{0.8em}
\textbf{[Output Format]} \\
Knowledge: \textless knowledge\_name\textgreater \\
Bucket: \textless bucket\_name\textgreater

\end{tcolorbox}
\caption{Structured Prompt Template for the Exercise Selection Agent in AgentCAT}
\label{selection_agent_prompt}
\end{figure}